\documentclass[tikz,border=8pt]{standalone}
\usepackage{ctex}
\usepackage{amsmath}
\usetikzlibrary{calc, arrows.meta, 3d}

\begin{document}
% 直线与平面的三种位置关系
% 使用等距投影模拟3D平面
\begin{tikzpicture}[
  x={(1cm,0cm)},
  y={(0.35cm,0.25cm)},
  z={(0cm,0.85cm)},
  thick,
  plane/.style={fill=blue!10, draw=blue!40, fill opacity=0.7},
  lline/.style={red, thick},
  dot/.style={fill=black, circle, inner sep=1.5pt},
]

%%% 栏1：直线在平面内 l ⊂ α
\begin{scope}[xshift=-6cm, yshift=0cm]
  \node[font=\small\bfseries, align=center] at (1.2, 4.2) {(1) 直线在平面内};
  \node[font=\small] at (1.2, 3.8) {$l \subset \alpha$};

  % 平面四边形（等距投影）
  \fill[plane] (0,0) -- (2.4,0) -- (3.2,1.6) -- (0.8,1.6) -- cycle;
  \draw[blue!50] (0,0) -- (2.4,0) -- (3.2,1.6) -- (0.8,1.6) -- cycle;
  \node[blue!70, font=\small] at (3.4, 0.8) {$\alpha$};

  % 直线在平面内（水平方向）
  \draw[lline] (0.3, 0.8) -- (2.9, 0.8);
  \node[red, font=\small, above] at (1.6, 0.8) {$l$};
\end{scope}

%%% 栏2：直线与平面平行 l ∥ α
\begin{scope}[xshift=0cm, yshift=0cm]
  \node[font=\small\bfseries, align=center] at (1.2, 4.2) {(2) 直线与平面平行};
  \node[font=\small] at (1.2, 3.8) {$l \parallel \alpha$};

  % 平面
  \fill[plane] (0,0) -- (2.4,0) -- (3.2,1.6) -- (0.8,1.6) -- cycle;
  \draw[blue!50] (0,0) -- (2.4,0) -- (3.2,1.6) -- (0.8,1.6) -- cycle;
  \node[blue!70, font=\small] at (3.4, 0.8) {$\alpha$};

  % 直线在平面上方平行（稍高）
  \draw[lline] (0.1, 2.8) -- (2.7, 2.8);
  \node[red, font=\small, above] at (1.4, 2.8) {$l$};

  % 虚线辅助（显示平行）
  \draw[gray, dashed, thin] (1.4, 0.8) -- (1.4, 2.8);
\end{scope}

%%% 栏3：直线与平面相交 l ∩ α = P
\begin{scope}[xshift=6cm, yshift=0cm]
  \node[font=\small\bfseries, align=center] at (1.2, 4.2) {(3) 直线与平面相交};
  \node[font=\small] at (1.2, 3.8) {$l \cap \alpha = P$};

  % 平面
  \fill[plane] (0,0) -- (2.4,0) -- (3.2,1.6) -- (0.8,1.6) -- cycle;
  \draw[blue!50] (0,0) -- (2.4,0) -- (3.2,1.6) -- (0.8,1.6) -- cycle;
  \node[blue!70, font=\small] at (3.4, 0.8) {$\alpha$};

  % 交点 P（在平面内偏中位置）
  \coordinate (P) at (1.6, 0.8);

  % 直线穿过平面：平面下方部分（虚线）
  \draw[lline, dashed] (P) -- (2.5, -0.6);
  % 平面上方部分（实线）
  \draw[lline] (P) -- (0.5, 3.2);
  \node[red, font=\small, above left] at (0.5, 3.2) {$l$};

  % 交点标注
  \node[dot] at (P) {};
  \node[font=\small, right] at ($(P)+(0.1,0)$) {$P$};
\end{scope}

\end{tikzpicture}
\end{document}
